\sectionpaged{Dew point}
The \texttt{Dew point}\index{Dew point} is the temperature the air needs to be cooled to at constant pressure in order
to produce a relative humidity of 100\percent.
This temperature depends on the pressure and water content of the air.
When the air is cooled below the dew point, its moisture capacity is reduced and
airborne water vapor will condense to form liquid water known as dew.
When this occurs through the air's contact with a colder surface, dew will form on that surface.

The dew point is calculated from the following formulae\sidesupercite{WikiDewPoint}:

\begin{tabular}{c c c}
    $\AmbientTemperature$ & \textbf{b} & \textbf{c} \\
    \hline
    $> 0$ & 17.368 & 238.88 \\
    $\leq 0$ & 17.996 & 247.15 \\
    \hline
\end{tabular}

$ d = \dfrac{ b \times \AmbientTemperature }{ c + \AmbientTemperature } $

\NewValueUnit{\DewPoint}{T_{dp}}{Dew point Temperature}%
$ \DewPoint = c \times \log{ \dfrac{\RelativeHumidity}{100 \times e^d}} $

where:
\begin{description}[leftmargin=3em,style=nextline]
    \item[$\DewPoint$] the \texttt{Dew point} in degrees Celsius
    \item[$\AmbientTemperature$] the \texttt{Ambient temperature}\index{Temperature!Ambient} in degrees Celsius
    \item[$\RelativeHumidity$] the \texttt{Relative Humidity}\index{Relative Humidity}
\end{description}
